\section{Materials and Methods}
\label{sec:methods}

\begin{figure*}[t]
\centering
\begin{tikzpicture}[
  node distance=4mm and 6mm,
  box/.style={draw, rounded corners=2pt, align=center, font=\scriptsize, minimum height=8mm, inner sep=3pt},
  data/.style={box, fill=blue!6},
  proc/.style={box, fill=green!7},
  ui/.style={box, fill=orange!10},
  arr/.style={-{Stealth[length=2mm]}, semithick}]
\node[data] (ds) {TrialGPT Criterion\\Annotations (1,015)};
\node[proc, right=of ds] (split) {drop GPT-4 columns\\patient split 37/8/8\\(rng 42) + manifest};
\node[proc, right=of split] (prompt) {numbered note +\\criterion $\rightarrow$ chat prompt\\target: JSON label+IDs};
\node[proc, right=of prompt] (lora) {LoRA fine-tuning\\Qwen2.5-1.5B (T4)\\val selection};
\node[data, right=of lora] (adapter) {adapter\\(HF Hub)};
\node[proc, below=9mm of prompt] (dec) {own Qwen2 forward\\explicit KV cache\\+ prefix reuse};
\node[proc, left=of dec] (attn) {attention backend:\\SDPA / reference /\\tiled online-softmax};
\node[proc, right=of dec] (spec) {greedy speculative\\decoding (0.5B draft)\\cache rollback};
\node[proc, right=of spec] (post) {parse JSON, verify IDs,\\evidence-only explanation,\\rule aggregation};
\node[ui, left=of attn] (ui) {React UI\\FastAPI server\\(Docker)};
\draw[arr] (ds) -- (split);
\draw[arr] (split) -- (prompt);
\draw[arr] (prompt) -- (lora);
\draw[arr] (lora) -- (adapter);
\draw[arr] (adapter) |- ($(spec.north)+(0,3mm)$) -- (spec.north);
\draw[arr] (attn) -- (dec);
\draw[arr] (dec) -- (spec);
\draw[arr] (spec) -- (post);
\draw[arr] (ui) -- node[above, font=\tiny]{note, criteria} (attn);
\draw[arr] (post.south) -- ++(0,-4mm) -| node[pos=0.25, below, font=\tiny]{streamed results, report, benchmarks} (ui.south);
\end{tikzpicture}
\caption{System overview. Top: data preparation and supervised adaptation. Bottom: the inference
stack used by the application, built on our own decoder, attention and decoding implementations.}
\label{fig:overview}
\end{figure*}

Fig.~\ref{fig:overview} summarizes the system: the upper row prepares the data and adapts the model,
the lower row is the inference stack used for evaluation benchmarks and by the application.

\subsection{Dataset and split}
\label{sec:data}
The TrialGPT Criterion Annotations~\cite{jin2024trialgpt} contain 1,015 patient--criterion pairs
from 53 SIGIR-2016 patients~\cite{koopman2016test} and 103 trials. Each record holds the patient
note, already split into numbered sentences (3--14 per note, at most 235 words), the criterion type
and text, the expert eligibility label and the expert evidence sentence IDs. The record also stores
GPT-4's label, explanation and cited sentences, a judgement of that explanation
(\code{explanation\_correctness}) and the original authors' own \code{training} flag. All five fields
are removed when the dataset is loaded, so they cannot reach a prompt, a target or a metric.

The split follows the assignment protocol exactly: the 53 unique patient IDs are sorted,
shuffled with \code{numpy.random.default\_rng(42)}, and the first 37 go to training, the next 8 to
validation and the last 8 to testing. Every annotation of a patient therefore lives in exactly one
split; the preparation script asserts disjoint patient sets, single assignment of every annotation
and absence of leaky fields, and writes a manifest with patient IDs and label counts.
Table~\ref{tab:split} shows the resulting distribution. The label distribution is heavily skewed:
\emph{not excluded} accounts for almost half of all records, while \emph{excluded} occurs only 10
times overall and \emph{not at all} in the test split. One training annotation (ID 883) has an empty
criterion text; it remains assigned to its patient's split in the manifest but is excluded from
training because no task can be posed for it.

\begin{table}[t]
\centering
\caption{Patient-grouped split (seed 42) and label distribution.}
\label{tab:split}
\setlength{\tabcolsep}{4pt}
\begin{tabular}{llrrr}
\toprule
 & & Train & Val & Test \\
\midrule
\multicolumn{2}{l}{Patients} & 37 & 8 & 8 \\
\multicolumn{2}{l}{Annotations (used)} & 602 & 206 & 206 \\
\multicolumn{2}{l}{\quad with gold evidence} & 254 & 90 & 61 \\
\midrule
\multirow{4}{*}{Inclusion} & included & 88 & 35 & 27 \\
 & not included & 13 & 3 & 7 \\
 & not enough information & 109 & 27 & 38 \\
 & not applicable & 5 & 1 & 3 \\
\midrule
\multirow{4}{*}{Exclusion} & excluded & 8 & 2 & 0 \\
 & not excluded & 291 & 101 & 91 \\
 & not enough information & 65 & 26 & 28 \\
 & not applicable & 23 & 11 & 12 \\
\bottomrule
\end{tabular}
\end{table}

\subsection{Input/output format and tokenization}
\label{sec:format}
Each example becomes a two-turn chat in the model's template: a system message that frames the task
and states that missing information is never proof of eligibility or ineligibility, and a user
message with the numbered note, the criterion type, the criterion text, a one-line definition of each
of the four labels for that criterion type, and the required output schema. The target is a single
JSON object, e.g.\ \code{\{"label": "included", "evidence": [0, 3]\}}, followed by the end-of-turn
token. The four label sets are kept separate for inclusion and exclusion criteria, so the model always
chooses among the labels that are valid for the given type.

With the Qwen2.5 tokenizer (byte-level BPE, 151,665 tokens) prompts are 256--625 tokens long and
targets 12--38 tokens; the longest prompt plus target is 651 tokens. We therefore set the maximum
sequence length to 768, which means \emph{no example is truncated}. The training code raises an error
instead of truncating, and the application rejects (with a clear message) any input above 2,048
tokens, because silently cutting the note could remove exactly the evidence a decision depends on.

Generated text is parsed by extracting the first JSON object; a label outside the type's label set is
invalid, and cited IDs that do not exist in the note are discarded. If the JSON is malformed but a
valid label string appears, the label is recovered and the record is flagged.

\subsection{Model selection}
\label{sec:model}
Table~\ref{tab:models} compares the candidates we considered. All are decoder-only Transformers with
RoPE, RMSNorm, SwiGLU and GQA, so all support the required KV-cache and speculative-decoding work. Two
constraints dominated the choice: LoRA training must fit one free 16\,GB T4 GPU, and the application
must run on a CPU-only laptop with 16\,GB RAM. Speculative decoding additionally requires a smaller
model with an \emph{identical tokenizer}.

\begin{table}[t]
\centering
\caption{Candidate open-weight decoders (values from the model configurations).}
\label{tab:models}
\setlength{\tabcolsep}{2.5pt}
\begin{tabular}{lcccc}
\toprule
 & Qwen2.5 & Qwen2.5 & Llama 3.2 & Llama 3.2 \\
 & 1.5B-Instr. & 0.5B-Instr. & 1B-Instr. & 3B-Instr. \\
\midrule
Parameters & 1.54B & 0.49B & 1.24B & 3.21B \\
Layers / hidden & 28 / 1536 & 24 / 896 & 16 / 2048 & 28 / 3072 \\
Q / KV heads & 12 / 2 & 14 / 2 & 32 / 8 & 24 / 8 \\
Head dim & 128 & 64 & 64 & 128 \\
Vocabulary & 151,936 & 151,936 & 128,256 & 128,256 \\
Context & 32K & 32K & 128K & 128K \\
KV cache / token (fp16) & 28\,KB & 12\,KB & 32\,KB & 112\,KB \\
Weights fp16 / fp32 & 3.1 / 6.2\,GB & 1.0 / 2.0\,GB & 2.5 / 5.0\,GB & 6.4 / 12.8\,GB \\
License & Apache-2.0 & Apache-2.0 & Llama 3.2 & Llama 3.2 \\
\bottomrule
\end{tabular}
\end{table}

We selected \textbf{Qwen2.5-1.5B-Instruct} as the target and \textbf{Qwen2.5-0.5B-Instruct} as the
draft. The pair shares one tokenizer and chat template, is Apache-2.0 licensed with ungated downloads,
and the target's aggressive GQA (12 query heads, 2 KV heads) gives the smallest KV cache per token of
the larger candidates. Its fp32 weights (6.2\,GB) fit the laptop's RAM for CPU inference, whereas
Llama~3.2-3B (12.8\,GB in fp32) does not leave room for the draft, the cache and the OS. The
Llama~3.2-1B/3B pair was a viable alternative; we rejected it because of the gated license, the
larger KV cache of the 3B model, and because a 1B target would leave too little room between target
and draft for speculative decoding to pay off.

\subsection{Supervised fine-tuning}
\label{sec:finetune}
We adapt all seven linear projections of every layer ($W_q, W_k, W_v, W_o$ and the gate, up and down
projections of the MLP) with LoRA~\cite{hu2022lora}: $W' = W + \frac{\alpha}{r}BA$ with rank $r=16$,
$\alpha=32$ and dropout 0.05, i.e.\ 18.5M trainable parameters (1.18\%). The objective is the causal
language-modelling cross-entropy on the \emph{answer tokens only}; prompt positions are masked with
$-100$ so the model is not trained to reproduce notes or criteria. The frozen base weights are loaded
in fp16 (the T4 has no native bf16), adapters are kept in fp32, and training uses automatic mixed
precision with dynamic loss scaling.

The training loop is written in plain PyTorch rather than with a trainer abstraction: AdamW,
learning rate $2\cdot10^{-4}$ with 5\% linear warm-up and cosine decay, gradient clipping at 1.0,
micro-batch 2 with 4 accumulation steps (effective batch 8), 3 epochs, seed 42. A micro-batch of 4
ran out of GPU memory: with a 151,936-word vocabulary the fp32 logits of a single micro-batch of four
651-token sequences occupy $4\cdot651\cdot151{,}936\cdot4$\,B $\approx$ 1.6\,GB, and several copies
exist during the loss and backward computation. Gradient checkpointing keeps the peak at 5.5\,GB.
After every epoch the adapter is saved, the teacher-forced validation loss is computed, and the
validation split is decoded greedily and scored; the epoch with the highest validation macro-F1 is
kept. The test split is read only by the final evaluation command (Section~\ref{sec:setup}).

\textbf{Ablations.} Three alternatives were trained with otherwise identical settings and compared on
validation only: (i) \emph{label-first} answer order versus the \emph{evidence-first} order of our
initial design, which asks the model to commit to evidence before the label; (ii) \emph{rare-class
oversampling}, which repeats every training example $\min(8,\operatorname{round}(\sqrt{n_{\max}/n_c}))$
times per epoch, where $n_c$ is the count of its (criterion type, label) class (602 $\rightarrow$
1,060 examples per epoch); (iii) training for 5 instead of 3 epochs. The 0.5B draft model is
fine-tuned with the same recipe and the selected answer order.

\subsection{Explicit KV-cached decoding}
\label{sec:kv}
To control every step of cached decoding we implemented the Qwen2 forward pass ourselves on top of
the trained weight tensors (embedding, projections, norms); only the linear layers are reused as
modules. Per layer: $x \leftarrow h + W_o\,\mathrm{Attn}(\mathrm{RoPE}(W_q\bar h + b_q),
\mathrm{RoPE}(W_k\bar h + b_k), W_v\bar h + b_v)$ with $\bar h = \mathrm{RMSNorm}(h)$, followed by
$h \leftarrow x + W_d(\mathrm{SiLU}(W_g\bar x)\odot W_u\bar x)$.

\textbf{Cache layout and growth.} The cache holds, for each of the 28 layers, a key and a value
tensor of shape $(B, H_{kv}{=}2, L_{\max}, 128)$ that is \emph{pre-allocated} for the prompt length
plus the generation budget; a single integer \code{length} marks the valid prefix. A forward call on
$T$ new tokens at absolute position $p$ writes their keys and values into slots $p\ldots p{+}T{-}1$
and attends over slots $0\ldots p{+}T{-}1$. Growth therefore costs $O(1)$ per token instead of the
$O(n)$ reallocation of repeated concatenation, and rollback (needed by speculative decoding and by
prefix reuse) is just lowering \code{length}. Memory grows linearly: $2 \cdot 28 \cdot 2 \cdot 128$
values, i.e.\ 28\,KB per token in fp16. Keys and values are stored with the 2 KV heads and expanded to
the 12 query heads only inside attention, so GQA's 6$\times$ saving is kept.

\textbf{Positional handling.} RoPE rotates queries and keys by angles $\theta_i\,m$ that depend on
the absolute position $m$; cached keys were rotated when they were written. New tokens must therefore
be rotated with $m = p, \ldots, p{+}T{-}1$ rather than $0\ldots T{-}1$, otherwise their relative
offsets to the cached keys would be wrong. Our forward pass takes $p$ explicitly and computes
$\cos/\sin$ tables for exactly those positions (base $\theta = 10^6$).

\textbf{Masking.} With $T$ new queries and $p+T$ keys, query $i$ may attend key $j$ iff
$j \le i + p$: a causal mask aligned to the \emph{bottom-right} of the $T\times(p{+}T)$ score matrix.
PyTorch's \code{is\_causal} flag aligns to the top-left, which is only correct when $p=0$; we pass the
bottom-right mask explicitly when $p>0$ and $T>1$. For single-token decode steps ($T=1$) no mask is
needed because every cached position is in the past.

\textbf{Prefix reuse.} In the application all criteria of one patient share the system prompt and
the note. The cache remembers which token IDs its valid slots belong to; a new prompt reuses the
longest common token prefix (keeping at least one token to produce logits), crops the rest, and
prefills only the new suffix.

\textbf{Reference path and numerics.} The uncached reference allocates a fresh cache and runs the
full sequence at every step, i.e.\ $O(n^2)$ total work for $n$ tokens. In exact arithmetic both paths
compute the same logits. In floating point they need not be bitwise identical: the cached path
computes the new token's projections as a $1\times d$ product while the reference computes them as one
row of a $T\times d$ product, and different matrix shapes select different kernels and summation
orders. Differences are of the order of float rounding and can flip a greedy argmax only when the top
two logits are within that margin, which we test for explicitly (Section~\ref{sec:res-kv}).

\subsection{Tiled attention with online softmax}
\label{sec:tiled}
The reference attention computes $S = QK^\top/\sqrt{d}$, applies the mask with $-\infty$, subtracts
the row maximum before exponentiation and normalizes; it materializes the full $L_q\times L_k$ score
matrix. The tiled variant (Algorithm~\ref{alg:tiled}) processes queries in blocks of $B_q$ rows and
streams over key/value blocks of $B_k$ columns, keeping per row a running maximum $m$, a running
normalizer $\ell$ and an unnormalized output accumulator $o$~\cite{milakov2018online,dao2022flashattention}.
When a new block raises the maximum, previous contributions are rescaled by $e^{m-m'}$, so the result
is exactly softmax-weighted, while only a $B_q\times B_k$ tile of scores exists at any time. Under
causal masking, key blocks entirely above the diagonal are skipped. Half-precision inputs are
accumulated in fp32. Rows whose blocks are fully masked keep $m=-\infty$, so the subtraction uses a
safe maximum of 0 to avoid $-\infty-(-\infty)$.

\begin{algorithm}[t]
\caption{Tiled attention with online softmax (one query block)}
\label{alg:tiled}
\begin{algorithmic}[1]
\Require $Q_b\in\mathbb{R}^{B_q\times d}$, $K, V\in\mathbb{R}^{L_k\times d}$, scale $s$, mask
\State $m \gets -\infty^{B_q}$, $\ell \gets 0^{B_q}$, $o \gets 0^{B_q\times d}$
\For{each key block $j$ (skipping blocks above the causal diagonal)}
  \State $S \gets s\,Q_b K_j^\top$; set masked entries of $S$ to $-\infty$
  \State $m' \gets \max(m, \operatorname{rowmax}(S))$
  \State $P \gets \exp(S - m')$; \quad $c \gets \exp(m - m')$
  \State $\ell \gets c\odot\ell + \operatorname{rowsum}(P)$; \quad $o \gets c\odot o + P V_j$
  \State $m \gets m'$
\EndFor
\State \Return $o / \ell$
\end{algorithmic}
\end{algorithm}

Both kernels are registered as attention backends of the Hugging Face model
(\code{attn\_implementation="tiled"}), and our decoder can select them. Registering only the attention
function is not enough: the library then passes \emph{no} attention mask to an unknown backend, so a
left-padded batch silently attends to padding tokens (we measured logit errors up to 10.6). Registering
the SDPA mask builder under the same name fixes this.

\subsection{Greedy speculative decoding}
\label{sec:spec}
Algorithm~\ref{alg:spec} gives our greedy variant of speculative decoding~\cite{leviathan2023fast,chen2023accelerating}.
Both models keep their own KV cache with the invariant that, between iterations, each cache holds
every committed token except the last one, which is the pending input. In an iteration the draft first
catches up on committed tokens it has not processed, then proposes $k$ tokens greedily. The target
processes the last committed token and the $k$ drafts in \emph{one} forward pass; row $i$ of its
logits is the target's greedy choice after drafts $1\ldots i$. The longest prefix of drafts matching
the target's choices is accepted, and the target's own token at the first mismatch (or a bonus token
if all match) is appended, so every iteration commits between 1 and $k{+}1$ tokens. Finally both
caches are cropped to the committed length minus one, discarding the keys and values of rejected
drafts. Because every committed token equals the target's argmax given the committed prefix, the
output is identical to plain greedy decoding of the target.

\begin{algorithm}[t]
\caption{Greedy speculative decoding}
\label{alg:spec}
\begin{algorithmic}[1]
\Require prompt $x$, target $M_t$, draft $M_d$, draft length $k$
\State prefill $M_t$ on $x$; $y \gets x \,\|\, \arg\max M_t(x)$
\While{budget left and last token $\neq$ EOS}
  \State $M_d$: feed $y[\text{len}(C_d):]$, then draft $d_1..d_k$ greedily
  \State $M_t$: one pass on $[y_{-1}, d_1..d_k]$ $\rightarrow$ choices $g_0..g_k$
  \State $n \gets$ largest $n$ with $d_i = g_{i-1}\;\forall i\le n$
  \State $y \gets y \,\|\, d_1..d_n \,\|\, g_n$ \Comment{truncate at EOS / budget}
  \State crop $C_t$ and $C_d$ to $|y| - 1$
\EndWhile
\end{algorithmic}
\end{algorithm}

\subsection{Explanations, aggregation and application}
\label{sec:app}
The dataset contains no expert explanations, so the model is not trained to write free text.
Instead, each prediction's explanation is assembled \emph{only} from the cited sentences (quoted with
their IDs) and the meaning of the predicted label; a decisive label without cited evidence is flagged
\emph{unsupported}, and \emph{not enough information} lists the criterion as missing information.
This avoids hallucinated rationales at the cost of fluency. The overall patient--trial outcome is
an explicit rule, not a learned or ground-truth label: \emph{likely ineligible} if any inclusion
criterion is \emph{not included} or any exclusion criterion is \emph{excluded}; otherwise \emph{needs
clinician review} if any criterion has insufficient information or could not be parsed; otherwise
\emph{potentially eligible}.

The application (Fig.~\ref{fig:overview}, bottom) is a FastAPI server with a React front end,
packaged as a Docker image and started with one command. The user pastes a note (or loads a
validation example), enters or pastes criteria or fetches them from ClinicalTrials.gov by NCT ID,
chooses a decoding mode and runs the screening. Results stream back criterion by criterion with the
label, evidence sentences highlighted in the numbered note, explanation, flags and timing. Reports are
saved with provenance (model, adapter, decoding settings, time) and can be downloaded as HTML or JSON.
A benchmark page runs a chosen prompt live through the uncached, cached, tiled-attention and
speculative decoders and shows timings, token equality and draft acceptance, and an experiments page
renders training curves, confusion matrices and GPU benchmarks from the logged result files.
